\clearpage
\section{계산 절차와 장 요약}
\label{sec:quartic-quintic-workflow}

\begin{learninggoal}
사차·오차 갈루아군 계산을 재현 가능한 단계로 정리한다. 각 단계가 제공하는 정보와 제공하지 못하는 정보를 구분한다.
\end{learninggoal}

\subsection{기약 사차식 계산 절차}

$K$의 표수가 $2,3$이 아니고 $f\in K[X]$가 일계수 사차다항식이라 하자.

\begin{enumerate}
  \item \textbf{분리성과 기약성 확인}: $\gcd(f,f')=1$인지 확인하고, 기약이면 갈루아군이 추이적임을 기록한다.
  \item \textbf{삼차항 제거}: $X=Y-a/4$를 대입하여
  \[
    Y^4+pY^2+qY+r
  \]
  꼴로 바꾼다. 평행이동은 분해체와 판별식을 바꾸지 않는다.
  \item \textbf{삼차분해식 계산}:
  \[
    R_f(Z)=Z^3-2pZ^2+(p^2-4r)Z+q^2.
  \]
  \item \textbf{인수분해형 판정}:
  \[
  \begin{array}{ccl}
    1+1+1&\Rightarrow&V_4,\\
    3&\Rightarrow&A_4\text{ 또는 }S_4,\\
    1+2&\Rightarrow&C_4\text{ 또는 }D_4.
  \end{array}
  \]
  \item \textbf{판별식 또는 차수로 마무리}: 기약 분해식이면 판별식 제곱 여부로 $A_4/S_4$를 가른다. $1+2$형이면 분해식 분해체 $F$ 위의 제곱근 첨가 차수 $[L:F]$ 또는 명제 \ref{prop:c4-d4-square-test}를 사용한다.
  \item \textbf{독립 검산}: 가능한 경우 좋은 소수의 인수분해형, 분해체의 직접 구성, 실제 근의 형태 중 하나로 다시 확인한다.
\end{enumerate}

\begin{keyidea}
사차식에서는 하나의 삼차분해식이 거의 모든 군론적 정보를 압축한다. 그 인수분해형은 $G/(G\cap V_4)$을, 판별식은 $G\cap A_4$를, 마지막 제곱근 차수는 $G\cap V_4$의 크기를 측정한다.
\end{keyidea}

\subsection{기약 오차식 계산 절차}

$f\in\Z[X]$가 일계수 기약 오차다항식이라 하자.

\begin{enumerate}
  \item \textbf{분리성·기약성·판별식}: 기약성을 먼저 증명하고 $\Disc(f)$의 제곱 여부로
  \[
    \{C_5,D_5,A_5\}
    \quad\text{또는}\quad
    \{F_{20},S_5\}
  \]
  중 어느 쪽인지 정한다.
  \item \textbf{좋은 소수 선택}: 판별식을 나누지 않는 작은 소수들을 택하고 $\bar f$를 정확히 인수분해한다.
  \item \textbf{순환형 누적}: 서로 다른 좋은 소수에서 얻은 순환형들을 후보군 표와 대조한다. 전치나 $3+2$형을 찾으면 즉시 $S_5$이다. 삼순환과 제곱 판별식을 찾으면 $A_5$이다.
  \item \textbf{작은 군의 분해식}: 후보가 $C_5/D_5$로 남으면 $D_5$-분해식을, $F_{20}/S_5$로 남는데 강한 순환형이 없으면 적절한 부분군 분해식을 사용한다.
  \item \textbf{구조적 예외 활용}: $X^5-a$처럼 근의 형태가 명확하면 단위근을 직접 첨가하여 분해체 차수와 반직접곱을 계산한다. 실근 개수를 알 수 있으면 복소켤레의 순환형을 이용한다.
  \item \textbf{증명서 남기기}: 사용한 소수, 법 $p$ 인수분해, 각 인수의 기약성, 판별식과 소수의 서로소성, 분해식의 정확한 인수분해를 모두 기록한다.
\end{enumerate}

\subsection{각 도구의 한계}

\begin{center}
\renewcommand{\arraystretch}{1.22}
\begin{tabular}{p{0.23\textwidth}p{0.31\textwidth}p{0.34\textwidth}}
\toprule
도구 & 알려 주는 것 & 혼자서는 모르는 것 \\
\midrule
기약성 & 근 작용의 추이성 & 어느 추이적 부분군인지 \\
판별식 & $G\subseteq A_n$ 여부 & $A_n$ 안의 세부 부분군 \\
좋은 소수 & 특정 순환형 원소의 존재 & 군의 모든 원소와 전체 위수 \\
삼차분해식 & 사차군의 $S_3$ 몫작용 & $C_4/D_4$의 마지막 구분 \\
부분군 분해식 & 특정 부분군의 켤레 안에 포함되는지 & 값 충돌 시의 정확한 포함관계 \\
분해체 차수 & 갈루아군의 위수 & 같은 위수의 비동형군 구분 \\
\bottomrule
\end{tabular}
\end{center}

\begin{remark}[원자료와 이 장의 보강 범위]
원자료의 직접적인 중심은 사차식에서 세 쌍분할로 얻는 삼차분해식, $V_4$ 고정체와 사차 근호공식이다. 또한 기약 소수차수 방정식의 가해성 논의에서 $X^5-4X+2$의 $S_5$ 예제를 제시한다. 이 장의 추이적 부분군 완전분류, 좋은 소수 순환형 알고리즘, 일반 부분군 분해식과 $C_5,D_5,F_{20},A_5,S_5$의 다섯 계산 예제는 실제 계산을 위해 별도로 조직한 학습용 보강이다. 특히 데데킨트--프로베니우스 정리의 완전한 증명은 이 교재의 현재 선수범위를 넘어선다.
\end{remark}

\begin{chapterreview}
\begin{itemize}
  \item 네 근을 두 쌍으로 나누는 세 방법에 대한 $S_4$의 작용은 전사준동형 $S_4\to S_3$을 주며 핵은 $V_4$이다.
  \item $X^4+pX^2+qX+r$의 삼차분해식은
  \[
    Z^3-2pZ^2+(p^2-4r)Z+q^2
  \]
  이고 원래 사차식과 판별식이 같다.
  \item 기약 사차식의 추이적 갈루아군은 $V_4,C_4,D_4,A_4,S_4$뿐이다.
  \item 삼차분해식의 인수분해형과 판별식, 분해체 차수를 결합하면 다섯 군을 모두 구분한다.
  \item 좋은 소수에서 법 $p$ 기약인수의 차수는 갈루아군의 프로베니우스 순환형을 준다.
  \item 기약 오차식의 추이적 갈루아군은 $C_5,D_5,F_{20},A_5,S_5$뿐이다.
  \item 제곱 판별식은 후보를 $C_5,D_5,A_5$로, 비제곱 판별식은 $F_{20},S_5$로 줄인다.
  \item $D_5$-분해식은 오각형의 변의 합 $\Theta_D$를 사용하며 차수는 $12$이다.
  \item 한 부분군 분해식의 유리근은, 값 충돌이 없을 때 갈루아군이 해당 부분군의 켤레 안에 들어감을 뜻한다.
  \item $X^5-4X+2$는 복소켤레의 전치 또는 법 $7$의 $3+2$형으로 $S_5$임을 판정할 수 있다.
\end{itemize}
\end{chapterreview}

\subsection*{복습 카드}

\begin{itemize}
  \item \textbf{$S_4\to S_3$의 핵은?} $V_4$.
  \item \textbf{사차 삼차분해식은?} $Z^3-2pZ^2+(p^2-4r)Z+q^2$.
  \item \textbf{기약 사차식의 다섯 군은?} $V_4,C_4,D_4,A_4,S_4$.
  \item \textbf{기약 오차식의 다섯 군은?} $C_5,D_5,F_{20},A_5,S_5$.
  \item \textbf{좋은 소수란?} $p\nmid\Disc(f)$인 소수.
  \item \textbf{법 $p$ 인수차수의 의미는?} 갈루아군 프로베니우스 원소의 순환길이.
  \item \textbf{$D_5$-분해식 차수는?} $[S_5:D_5]=12$.
  \item \textbf{오차식에서 전치가 발견되면?} 갈루아군은 $S_5$.
\end{itemize}
